\section{A label budget on horizontal lines}\label{sec:budget}

Fix a chart $I$ from \Cref{lem:charts}, with horizontal order $r$ and
cone slope $s>0$ as in \Cref{prop:cone}. For each deterministic label
$\gamma$ used below, $B_\gamma$ denotes the common Busemann function
of its coalescing family in the base environment.

The comparison concerns a bigeodesic in a finitely increased metric
whose upper label lies between two deterministic labels $\alpha,\beta$
with $r(\alpha)<r(\beta)$.
Suppose the increase adds a total cost $b$ on finite initial portions of
its arms from a site $k$. Under the cone bounds stated below, the base
upper $\alpha$-ray and lower $-\beta$-ray from a sufficiently distant
site $y$ to the right supply a bypass in unchanged weights. We will prove
that
\[
 b\le B_\alpha(y,k)+B_{-\beta}(y,k).
\]
For successive sites whose labels lie in the same interval, these
increments telescope. Their two linear parts leave the width
$r(\beta)-r(\alpha)$ as the available cost per horizontal unit.
We first use this budget in a one-edge resampling estimate needed to
exclude forks, and later to contradict the growing gains constructed in
\Cref{sec:tilts}.

Choose deterministic nested finite sets $D_1\subset D_2\subset\cdots$
of labels in $I$, containing both chart endpoints, such that their union
$D_I$ is dense in $I$ and $r(D_I)$ contains the endpoints of every
complementary interval of $r(I)$ in its convex hull. There are only
countably many such intervals, so these choices are possible. The bins
of $\Pi_j$ are
\[
 (\alpha,\beta)_I
   =\{\rho\in I:r(\alpha)<r(\rho)<r(\beta)\},
\]
where $\alpha,\beta$ are consecutive members of $D_j$ in $r$-order.
Empty bins may be retained. Their widths always sum to
\begin{equation}\label{eq:chart-width}
 W_I=r(\rho_r)-r(\rho_\ell).
\end{equation}
Every two distinct labels of $I\setminus D_I$ eventually belong to
different bins. If a label lies strictly between them, density supplies
an intervening member of $D_I$; otherwise they are endpoints of a
complementary interval and already belong to $D_I$.

By \Cref{cor:fixed-label}, almost surely no bigeodesic has an end label
in the deterministic countable set $D_I\cup(-D_I)$. The same exclusion
and the fixed-functional conclusions of \Cref{thm:structure} hold
simultaneously for all partition endpoints in all charts, and in every
countable collection of environments whose laws are absolutely
continuous with respect to the iid law. We work on the corresponding
common event. Every fixed-label comparator ray is calibrated by its
$B_\gamma$, by \Cref{lem:busemann-basic}.

\begin{lemma}[Comparison after a finite increase]\label{lem:budget-comparison}
Let $T$ be a base metric and $T^k$ a metric obtained by increasing finitely
many edge weights.  Both environments are assumed to have the structural
properties already established.  Suppose that $k,y$ lie on the same
horizontal line, with $k_1<y_1$, and that the following conditions hold for
a bin $(\alpha,\beta)_I$.
\begin{enumerate}[label=\textup{(\roman*)}]
\item Every changed edge has both endpoints in
$\{z:|z-k|_\infty\le R\}$, where $R\ge1$.
\item In $T^k$ there is an oriented bigeodesic through $k$ with upper
label $\rho\in(\alpha,\beta)_I$.  Its arms $p^+$ and $p^-$ are contained
in $Q^+(k,H_p)$ and $Q^-(k,H_p)$, respectively.
\item Finite initial portions of these two arms have all their vertices
in the same ball of radius $R$, and the sum of the increases along those
portions is at least $b>0$.
\item The base rays from $y$ of labels $\alpha$ and $-\beta$ are contained
in $Q^+(y,H_c)$ and $Q^-(y,H_c)$, respectively.
\end{enumerate}
There is a separation $S=S(s,R,H_p,H_c)$ such that $y_1-k_1>S$ implies
\begin{equation}\label{eq:budget-comparison}
 b\le B_\alpha(y,k)+B_{-\beta}(y,k).
\end{equation}
One may take
\begin{equation}\label{eq:budget-separation}
 S=\max\left\{R+\frac{R+H_c}{s},\frac{H_p+H_c}{s}\right\}
 \le C_s(1+R+H_p+H_c),
\end{equation}
with $C_s$ depending only on $s$.  In particular, the choice
of $S$ is uniform over bins and does not depend on the magnitudes of the
individual increases.
\end{lemma}

\begin{proof}
Put $d=y_1-k_1$.  If a vertex $w$ belongs to the ball of radius $R$ about
$k$, then
\[
 |w_2-y_2|\le R,
 \qquad |w_1-y_1|\ge d-R.
\]
Membership in either comparator cone would therefore require
$R\ge s(d-R)-H_c$.  Thus
\begin{equation}\label{eq:avoid-gain-ball}
 d>R+\frac{R+H_c}{s}
\end{equation}
makes both base comparator rays avoid the entire ball.  Their edge
weights are unchanged in $T^k$.  Since every other edge weight only
increases, each comparator remains a geodesic ray in $T^k$.
\Cref{lem:finite-change} preserves its label.

Increase the separation further, if necessary, to apply
\Cref{lem:crossing} with offsets $H_p,H_c$.  The upper marked ray is to
the left at its starting site but has label strictly east of $\alpha$.
The same ordering holds for the lower rays, since
\[
 r(-\beta)=-r(\beta)<-r(\rho)=r(-\rho).
\]
The upper comparator therefore meets $p^+$, and the lower comparator
meets $p^-$.  Call intersection vertices $w^+$ and $w^-$, respectively.
By \eqref{eq:avoid-gain-ball}, each intersection is outside the ball.
Consequently it occurs after the whole prescribed initial gain portion
of the corresponding marked arm.  This remains true if an arm leaves
and later reenters the ball, since every vertex of its prescribed initial
portion lies in the ball.

Let $g^+$ and $g^-$ be the increases on those portions, so
$g^++g^-\ge b$.  Writing $\operatorname{cost}_T$ for the cost of a
specified finite path in the base weights, domination by passage time and
then by a Busemann increment gives
\[
 \begin{split}
 \operatorname{cost}_{T^k}(p^+[k,w^+])
 &\ge g^++\operatorname{cost}_T(p^+[k,w^+])
 \ge g^++B_\alpha(k,w^+),\\
 \operatorname{cost}_{T^k}(p^-[k,w^-])
 &\ge g^-+\operatorname{cost}_T(p^-[k,w^-])
 \ge g^-+B_{-\beta}(k,w^-).
 \end{split}
\]
The marked segment from $w^-$ through $k$ to $w^+$ is minimizing in
$T^k$.  A competing finite walk follows the lower comparator backwards
from $w^-$ to $y$, then follows the upper comparator to $w^+$.
The cost of this walk in $T^k$ equals its base cost, namely
\[
 B_{-\beta}(y,w^-)+B_\alpha(y,w^+).
\]
No disjointness between the two comparator portions is required: a finite
walk is an admissible competitor and its cost counts every traversal.
Comparison with the marked segment yields
\[
 b+B_\alpha(k,w^+)+B_{-\beta}(k,w^-)
 \le B_\alpha(y,w^+)+B_{-\beta}(y,w^-).
\]
Additivity cancels the increments ending at the intersection vertices
and gives exactly \eqref{eq:budget-comparison}.  Both the avoidance bound
and the separation in \Cref{lem:crossing} have the asserted dependence
on the offsets and $R$.
\end{proof}

\Needspace{4\baselineskip}
\begin{proposition}[The label budget]\label{prop:label-budget}
Fix a bin $(\alpha,\beta)_I$, an integer spacing $l\ge1$, and a gain
$b>0$.  On a common probability space, mark some sites of the horizontal
grid $k_j=(jl,0)$, $j\in\Z$.  Suppose all sites have the same marginal
probability $q_{\alpha,\beta}$ of being marked in this bin, and the base
environment has its usual fixed-label Busemann asymptotics.  Suppose also
that for successive marks $k<y$ in the bin the comparison
\eqref{eq:budget-comparison} holds.  In particular, this is so if the
hypotheses of \Cref{lem:budget-comparison} hold with common deterministic
bounds and successive marks are sufficiently separated.  Then
\begin{equation}\label{eq:label-budget}
 \frac{b}{l}\,q_{\alpha,\beta}
 \le r(\beta)-r(\alpha).
\end{equation}
No independence or ergodicity assumption on the marks is needed.
\end{proposition}

\begin{proof}
Set $w=r(\beta)-r(\alpha)>0$ and, for $j\ge0$, write
\[
 A_j=B_\alpha(0,k_j)+B_{-\beta}(0,k_j)
     =-wjl+e_j.
\]
The sign in this formula uses
$\alpha(e_1)+(-\beta)(e_1)=r(\alpha)-r(\beta)=-w$.
The two fixed-label asymptotics give $e_j=o(jl)$ as $j\to\infty$.
Define
\[
 E_N=\max_{0\le j\le N}|e_j|.
\]
Then $E_N/(Nl)\to0$ almost surely.  To see the needed uniformity, for
each $\delta>0$ choose $J$ so that $|e_j|\le\delta jl$ when $j\ge J$.
The earlier errors have a finite maximum $C_J$, and
$E_N\le\max\{C_J,\delta Nl\}$.  Let $N\to\infty$ and then
$\delta\downarrow0$.  This argument controls random first and last marks
without replacing them by deterministic endpoints.

Let $m_N$ be the number of marks in $\{k_0,\ldots,k_N\}$.  If
$m_N\ge2$, list their indices as $i_1<\cdots<i_{m_N}$.  Summing the
comparison over consecutive marks and using additivity gives
\[
 \begin{split}
 b(m_N-1)
 &\le A_{i_1}-A_{i_{m_N}}\\
 &=wl(i_{m_N}-i_1)+e_{i_1}-e_{i_{m_N}}
 \le wNl+2E_N.
 \end{split}
\]
If $m_N=0$ or $1$, the same upper bound applies to
$b(m_N-1)_+$.  Therefore, in all cases,
\begin{equation}\label{eq:bounded-budget-count}
 bm_N\le b+wNl+2E_N,
 \qquad
 \limsup_{N\to\infty}\frac{m_N}{Nl}\le\frac wb
 \quad\text{almost surely}.
\end{equation}
The expectation step uses only the deterministic bounds
\[
 0\le\frac{m_N}{Nl}\le\frac{N+1}{Nl}\le\frac2l,
 \qquad N\ge1.
\]
Fatou's lemma applied to the nonnegative variables
$2/l-m_N/(Nl)$ gives the bounded reverse Fatou inequality
\[
 \limsup_{N\to\infty}\EE\frac{m_N}{Nl}
 \le\EE\left[\limsup_{N\to\infty}\frac{m_N}{Nl}\right]
 \le\frac wb.
\]
Equal marginal probabilities imply
$\EE m_N=(N+1)q_{\alpha,\beta}$ by linearity of expectation.  Hence
$q_{\alpha,\beta}/l\le w/b$, as asserted.  In particular, this proof
takes no expectation of either a Busemann error or a Busemann increment.
\end{proof}

We now prepare the one-edge resampling estimate. Replacing an edge by
an independent sample keeps the output environment iid. Choosing the old
weight below one level and the replacement above a larger level creates
a fixed positive gain. The same thresholds will ensure that long simple
paths contain many edges above the larger level. These edges supply the
edge--label credits in the fork count and the accumulated gains in the
later perturbations.

\begin{lemma}[Usable edges on long simple paths]\label{lem:usable}
There exist $a_0<b_0$ such that
\begin{equation}\label{eq:low-usable-probabilities}
 \begin{gathered}
 p_a=\PP(t_e\le a_0)>0,
 \qquad q=\PP(t_e<b_0)<\frac1{64},\\
 p_b=\PP(t_e\ge b_0)=1-q>0.
 \end{gathered}
\end{equation}
Call an edge \emph{usable} if $t_e\ge b_0$, and put
$\theta=8\sqrt q<1$.  For each vertex $x$ and integer $n_0\ge1$,
\begin{equation}\label{eq:usable-fixed-site}
 \PP\left(
 \begin{array}{c}
 \text{some simple path from $x$ of length $j\ge n_0$}\\
 \text{has fewer than $j/2$ usable edges}
 \end{array}\right)
 \le\frac{\theta^{n_0}}{1-\theta}.
\end{equation}
If $U_n$ is the event that every simple path starting in
$[-3n,3n]^2\cap\Z^2$ and having length at least $n/2$ has at least half
its edges usable, then
\begin{equation}\label{eq:usable-box}
 \PP(U_n^c)
 \le (6n+1)^2\frac{\theta^{\lceil n/2\rceil}}{1-\theta}
 \longrightarrow0.
\end{equation}
Almost surely, at every fixed vertex all sufficiently long simple paths
satisfy the same fraction bound.
\end{lemma}

\begin{proof}
Let $a$ be the bottom of the support of $G$.  It is finite, and
$G(\{a\})=0$.  As $b\downarrow a$ with $b>a$, the probability
$G([0,b))$ tends to zero and is positive.  Choose $b_0>a$ so that this
probability is less than $1/64$, and choose $a_0\in(a,b_0)$.
The definition of $a$ gives $p_a>0$, proving
\eqref{eq:low-usable-probabilities}.

The edges of a simple path are distinct.  For a specified path of length
$j$, independence and a union bound over subsets of its edges bound the
probability that at least $j/2$ of them are unusable by
\[
 2^j q^{j/2}.
\]
This also bounds the probability of fewer than $j/2$ usable edges.
There are at most $4^j$ nearest-neighbor paths of length $j$ from a given
vertex, so the union bound over simple paths is at most $\theta^j$.
Summing over $j\ge n_0$ proves \eqref{eq:usable-fixed-site}; the further
union bound over the $(6n+1)^2$ starting vertices proves
\eqref{eq:usable-box}.  Finally, the events in
\eqref{eq:usable-fixed-site} decrease as $n_0\to\infty$, and their
probabilities tend to zero.  Countability gives the last assertion
simultaneously at all vertices.  These estimates use no moment of an
individual edge weight.
\end{proof}

For an edge $e=\{x,x+v\}$, where $v$ is a lattice unit vector, and
$H\ge1$, let $\mathcal R_e(H)$ be the following regularity event in the
environment under consideration: from each endpoint $z$ of $e$, every
ray with label in $I$ is contained in $Q^+(z,H)$, and every ray with label
in $-I$ is contained in $Q^-(z,H)$.  Set $N_e(H)=0$ if $e$ is unusable
or $\mathcal R_e(H)$ fails.  Otherwise define
\begin{equation}\label{eq:edge-label-count-definition}
 N_e(H)=\#\{\rho\in I:\text{$e$ lies on a bigeodesic of upper label $\rho$}\}.
\end{equation}
Here the count is the number of distinct labels, and is assigned the
value $\infty$ if the set is infinite.  This convention does not
presuppose that the label set is countable.  By \Cref{lem:measurable},
existence of finitely many distinct witnessing labels is measurable in
the completed space.  The universal regularity event is imposed as a
separate environment-only filter; it is not inserted into an
existential path projection.

\begin{proposition}[Expected label count at an edge]\label{prop:edge-count}
There is a finite constant $C$, depending only on the chart and the weight
law, such that for every lattice edge $e$ and every $H\ge1$,
\begin{equation}\label{eq:edge-label-count}
 \EE N_e(H)\le C(H+1).
\end{equation}
\end{proposition}

\begin{proof}
Fix a unit vector $v$ and let $e_k=\{k,k+v\}$ for sites $k$ on a
horizontal grid whose spacing will be chosen below.  Use one iid base
environment $t$ and independent samples $\widehat t_{e_k}$ with law
$G$.  The individual environment $t^k$ equals $t$ off $e_k$ and equals
$\widehat t_{e_k}$ on $e_k$.  Each individual environment has precisely
the iid law.  We take the common full-measure event for the base and all
these countably many individual environments, including the fixed-label
exclusions at every endpoint in $D_I$ and $-D_I$.

For a bin $(\alpha,\beta)_I$, let $A_e^{\alpha,\beta}(H)$ denote the
event in an iid environment that $e$ is usable, $\mathcal R_e(H)$ holds,
and a bigeodesic whose upper label is in this bin uses $e$.  Mark $k$ in
this bin when
\begin{equation}\label{eq:one-edge-mark}
 t_{e_k}\le a_0
 \quad\text{and}\quad
 A_{e_k}^{\alpha,\beta}(H)\text{ holds in }t^k.
\end{equation}
The original variable $t_{e_k}$ is independent of the entire individual
environment $t^k$: that environment consists of the independent sample
$\widehat t_{e_k}$ and the other base coordinates.  Consequently the
common marginal marking probability is exactly
\begin{equation}\label{eq:one-edge-mark-probability}
 q_{\alpha,\beta}
   =p_a\,\PP\bigl(A_e^{\alpha,\beta}(H)\bigr),
\end{equation}
where $e$ is any fixed translate of $\{0,v\}$.
The regularity test is measurable in the individual environment, so it
does not alter this independence calculation.  Dependence between marks
at different sites is immaterial.

On a marking event, the one changed edge is increased by at least
$b_0-a_0$.  Every witnessing bigeodesic passes through $k$, and the edge
$e_k$ is the first edge of one of its arms based there.  Use that edge as
the gain portion of that arm and use the empty portion of the other arm.
Thus the modification and both gain portions fit in a ball of radius
one, and the total gain is at least
\[
 b=b_0-a_0>0.
\]
The regularity event in $t^k$ confines both marked arms with offset $H$.

We must also establish confinement of the \emph{base} comparator rays at
every marked site; base regularity was not assumed in
\eqref{eq:one-edge-mark}.  Let $q$ be any base ray from $k$ with label in
$I$ or $-I$.  A simple path from $k$ can traverse $e_k$ only as its first
edge.  Indeed, a later traversal in either direction would revisit $k$.
If $q$ avoids $e_k$, all its weights are unchanged and it remains a ray
in $t^k$ by monotonicity.  Its label is unchanged by
\Cref{lem:finite-change}, so $\mathcal R_{e_k}(H)$ in $t^k$ puts it in
the appropriate cone at $k$.

If $q$ uses $e_k$ first, its tail from $k+v$ avoids $e_k$ and remains a
ray in $t^k$, with the same label.  That tail lies in the corresponding
cone at $k+v$ with offset $H$.  For either choice of sign, the inequality
\[
 \pm(w_2-k_2-v_2)\ge s|w_1-k_1-v_1|-H
\]
implies
\[
 \pm(w_2-k_2)
 \ge s|w_1-k_1|-\bigl(H+|v_2|+s|v_1|\bigr).
\]
Including the initial vertex $k$, and then the initial edge by convexity,
confines the whole base ray with offset at most $H+1+s$.
This proves the needed base comparator control at both fixed labels
$\alpha$ and $-\beta$, uniformly over all bins.  It explains why
regularity was required at both endpoints of $e_k$.

Choose an integer grid spacing $l=l(H)$ large enough for
\Cref{lem:budget-comparison} with $R=1$, $H_p=H$, and $H_c=H+1+s$.
It may be chosen so that
\begin{equation}\label{eq:one-edge-spacing}
 2\le l(H)\le C_0(H+1)
\end{equation}
for a chart-dependent constant $C_0$, uniformly in the unit vector $v$.
The lower bound makes the edges $e_k$ at distinct grid sites distinct.
Any two different marked sites are now sufficiently separated.  Applying
\Cref{prop:label-budget} and
\eqref{eq:one-edge-mark-probability} to each bin yields
\[
 \PP\bigl(A_e^{\alpha,\beta}(H)\bigr)
 \le\frac{l(H)}{(b_0-a_0)p_a}
       \bigl(r(\beta)-r(\alpha)\bigr).
\]
For the finite partition $\Pi_j$, let $K_j(e,H)$ count its occupied bins
among the labels counted by $N_e(H)$.  Linearity of expectation and
\eqref{eq:chart-width} give
\begin{equation}\label{eq:occupied-bin-bound}
 \EE K_j(e,H)
 \le\frac{l(H)W_I}{(b_0-a_0)p_a}.
\end{equation}

On the event excluding $D_I$ as bigeodesic labels, refinement never loses a
counted label, so $K_j(e,H)$ is nondecreasing.  Every finite collection of
distinct counted labels is eventually in distinct bins.  Thus
\[
 K_j(e,H)\uparrow N_e(H),
\]
also when the latter is infinite: apply the separation assertion to
arbitrarily large finite collections.  Monotone convergence in
\eqref{eq:occupied-bin-bound}, followed by \eqref{eq:one-edge-spacing},
proves \eqref{eq:edge-label-count}.  Translation invariance and the
finite number of edge directions give one constant for all $e$.
\end{proof}
